\ifdefined\gkver\else\def\gkver{student}\fi
\documentclass[\gkver]{gaokaozhenti}
\begin{document}

\chapter{2008年上海}

\begin{enumerate}
\item
%% number: 1
%% paperName: 2008年上海高考第 1 题 4 分
%% typeId: 3
%% type: 填空题
%% chapter: 气体性质
%% point: 阿伏伽德罗常数
%% method: 无
%% score: 4
%% degree: 650
%% duplicateId: 0
%% body:
\textbf{A．}某行星绕太阳的运动可近似看作匀速圆周运动，已知行星运动的轨道半径为 $R$，周期为 $T$，万有引力恒量为 $G$，则该行星的线速度大小为 \tkanswer[2.6cm]{$\frac{2\pi R}{T}$}，太阳的质量可表示为 \tkanswer[2.8cm]{$\frac{4\pi^{2}R^{3}}{GT^{2}}$}。

\textbf{B．}体积为 $V$ 的油滴，滴在平静的水面上，扩展成面积为 $S$ 的单分子油膜，则该油滴的分子直径约为 \tkanswer[1.6cm]{$\frac{V}{S}$}。已知阿伏伽德罗常数为 $N_{A}$，油的摩尔质量为 $M$，则一个油分子的质量为 \tkanswer[1.6cm]{$\frac{M}{N_{A}}$}。

%% answer:
\jdanswer{
\textbf{A．}$\frac{2\pi R}{T}$，$\frac{4\pi^{2}R^{3}}{GT^{2}}$。

\textbf{B．}$\frac{V}{S}$，$\frac{M}{N_{A}}$。
}

%% memo:
\memoanswer{
\textbf{A．}该行星的线速度大小为 $v=\frac{2\pi R}{T}$；由万有引力定律得 $G\frac{Mm}{R^{2}}=\frac{mv^{2}}{R}$，解得太阳的质量 $M=\frac{4\pi^{2}R^{3}}{GT^{2}}$。

\textbf{B．}单分子油膜可视为横截面积为 $S$、高度为分子直径 $D$ 的长方体，则体积 $V=SD$，故分子直径约为 $D=\frac{V}{S}$；取 $1\Umol$ 油，含有 $N_{A}$ 个油分子，则一个油分子的质量为 $m=\frac{M}{N_{A}}$。
}

\item
%% number: 2
%% paperName: 2008年上海高考第 2 题 4 分
%% typeId: 3
%% type: 填空题
%% chapter: 原子和原子核
%% point: 半衰期
%% method: 无
%% score: 4
%% degree: 650
%% duplicateId: 0
%% body:
\textbf{A．}如图所示，把电量为 $-5\times10^{-9}\UC$ 的电荷，从电场中的 A 点移到 B 点，其电势能 \tkanswer[1.4cm]{增大}（选填“增大”、“减小”或“不变”）；若 A 点的电势 $U_{A}=15\UV$，B 点的电势 $U_{B}=10\UV$，则此过程中电场力做的功为 \tkanswer[2.6cm]{$-2.5\times10^{-8}\UJ$}。

\onepicture[w=4.0cm, align=right]{\includesvg{figs/02}}

\textbf{B．}放射性元素的原子核在 $\alpha$ 衰变或 $\beta$ 衰变生成新原子核时，往往会同时伴随着 \tkanswer[1.2cm]{$\gamma$} 辐射。已知 A、B 两种放射性元素的半衰期分别为 $T_{1}$ 和 $T_{2}$，$t=T_{1}\cdot T_{2}$ 时间后测得这两种放射性元素的质量相等，那么它们原来的质量之比 $m_{A}:m_{B}=$ \tkanswer[2.0cm]{$2^{T_{2}-T_{1}}$}。

%% answer:
\jdanswer{
\textbf{A．}增大，$-2.5\times10^{-8}\UJ$。

\textbf{B．}$\gamma$，$2^{T_{2}-T_{1}}$（或 $2^{T_{2}}:2^{T_{1}}$）。
}

%% memo:
\memoanswer{
\textbf{A．}将电荷从电场中的 A 点移到 B 点，电场力做负功，其电势能增加；由电势差公式 $U_{AB}=\frac{W}{q}$，得 $W=qU_{AB}=-5\times10^{-9}\times(15-10)\UJ=-2.5\times10^{-8}\UJ$。

\textbf{B．}放射性元素的原子核在 $\alpha$ 衰变或 $\beta$ 衰变生成新原子核时，往往以 $\gamma$ 光子的形式释放能量，即伴随 $\gamma$ 辐射；由半衰期的定义知，经过 $t=T_{1}\cdot T_{2}$ 时间后剩下的放射性元素的质量相同，则 $\frac{m_{A}}{2^{T_{2}}}=\frac{m_{B}}{2^{T_{1}}}$，故 $m_{A}:m_{B}=2^{T_{2}}:2^{T_{1}}$。
}

\item
%% number: 3
%% paperName: 2008年上海高考第 3 题 4 分
%% typeId: 3
%% type: 填空题
%% chapter: 交流电
%% point: 交流电
%% method: 无
%% score: 4
%% degree: 650
%% duplicateId: 0
%% body:
\textbf{A．}1911 年卢瑟福依据 $\alpha$ 粒子散射实验中 $\alpha$ 粒子发生了 \tkanswer[1.2cm]{大}（选填“大”或“小”）角度散射现象，提出了原子的核式结构模型。若用动能为 $1\UMeV$ 的 $\alpha$ 粒子轰击金箔，其速度约为 \tkanswer[2.2cm]{$6.9\times10^{6}$} $\Ums$。（质子和中子的质量均为 $1.67\times10^{-27}\Ukg$，$1\UMeV=10^{6}\UeV$）

\textbf{B．}某集装箱吊车的交流电动机输入电压为 $380\UV$，则该交流电压的最大值为 \tkanswer[1.6cm]{537} $\UV$。当吊车以 $0.1\Ums$ 的速度匀速吊起总质量为 $5.7\times10^{3}\Ukg$ 的集装箱时，测得电动机的电流为 $20\UA$，则电动机的工作效率为 \tkanswer[1.6cm]{$75\%$}。（$g$ 取 $10\Umsq$）

%% answer:
\jdanswer{
\textbf{A．}大，$6.9\times10^{6}\Ums$。

\textbf{B．}537，$75\%$。
}

%% memo:
\memoanswer{
\textbf{A．}卢瑟福在 $\alpha$ 粒子散射实验中发现了大多数 $\alpha$ 粒子没有大的偏转，少数发生了较大的偏转，卢瑟福抓住了这个现象进行分析，提出了原子的核式结构模型；由 $1\UMeV=1\times10^{6}\times1.6\times10^{-19}\UJ=\frac{1}{2}mv^{2}$，解得 $v=6.9\times10^{6}\Ums$。

\textbf{B．}输入电压 $380\UV$ 为有效值，则最大值为 $380\sqrt{2}\UV\approx537\UV$；电动机对集装箱做功的功率 $P=mgv=5.7\times10^{3}\times10\times0.1\UW=5.7\times10^{3}\UW$，电动机消耗的电功率 $P_{\text{总}}=380\times20\UW=7.6\times10^{3}\UW$，故电动机的工作效率为 $\eta=\frac{P}{P_{\text{总}}}=75\%$。
}

\item
%% number: 4
%% paperName: 2008年上海高考第 4 题 4 分
%% typeId: 3
%% type: 填空题
%% chapter: 运动和力
%% point: 牛顿第二定律
%% method: 图象法
%% score: 4
%% degree: 450
%% duplicateId: 0
%% body:
如图所示，在竖直平面内的直角坐标系中，一个质量为 $m$ 的质点在外力 $F$ 作用下，从坐标原点 O 由静止开始沿直线 ON 斜向下运动，直线 ON 与 $y$ 轴负方向成 $\theta$ 角（$\theta<\frac{\pi}{4}$）。则 $F$ 大小至少为 \tkanswer[1.8cm]{$mg\sin\theta$}；若 $F=mg\tan\theta$，则质点机械能大小的变化情况是 \tkanswer[2.6cm]{增大或减小}。

\onepicture[w=4.0cm, align=right]{\includesvg{figs/04}}

%% answer:
\jdanswer{$mg\sin\theta$，增大或减小}

%% memo:
\memoanswer{
该质点受到重力和外力 $F$ 从静止开始做直线运动，说明质点做匀加速直线运动；当 $F$ 的方向垂直于 ON 时，$F$ 最小为 $F=mg\sin\theta$。若 $F=mg\tan\theta$，则 $F$ 的方向可能与运动方向相同，也可能与运动方向相反，除重力外的 $F$ 对质点可能做正功，也可能做负功，故质点机械能增加、减少都有可能。

\onepicture[w=3.2cm]{figs/04_an.jpg}
}

\item
%% number: 5
%% paperName: 2008年上海高考第 5 题 4 分
%% typeId: 3
%% type: 填空题
%% chapter: 运动和力
%% point: 牛顿定律的应用
%% method: 无
%% score: 4
%% degree: 650
%% duplicateId: 0
%% body:
在伽利略羊皮纸手稿中发现的斜面实验数据如右表所示，人们推测第二、三列数据可能分别表示时间和长度。伽利略的一个长度单位相当于现在的 $\frac{29}{30}\Umm$，假设一个时间单位相当于现在的 $0.5\Us$。由此可以推算实验时光滑斜面的长度至少为 \tkanswer[1.6cm]{2.034} $\Um$；斜面的倾角约为 \tkanswer[1.4cm]{1.5} 度。（$g$ 取 $10\Umsq$）

\begin{table}[htbp]
\centering
\begin{tblr}{|c|c|c|}
\hline
\SetCell[c=3]{c} 表：伽利略手稿中的数据 & & \\
\hline
1 & 1 & 32 \\
\hline
4 & 2 & 130 \\
\hline
9 & 3 & 298 \\
\hline
16 & 4 & 526 \\
\hline
25 & 5 & 824 \\
\hline
36 & 6 & 1192 \\
\hline
49 & 7 & 1600 \\
\hline
64 & 8 & 2104 \\
\hline
\end{tblr}
\end{table}

%% answer:
\jdanswer{2.034，$1.5^{\circ}$}

%% memo:
\memoanswer{
依题意，第一列数据为时间的平方 $t^{2}$；从数据分析可知第一列数据与第三列数据之比约为 $1:32$（取平均值后比值为 $1:32.75$），即斜面长度与时间的平方成正比。由当时数据与现在的数据换算关系和匀变速运动公式，可得实验时光滑斜面的长度至少为 $2.034\Um$，斜面的倾角约为 $1.5^{\circ}$。
}

\item
%% number: 6
%% paperName: 2008年上海高考第 6 题 4 分
%% typeId: 1
%% type: 单选题
%% chapter: 原子和原子核
%% point: 人工核反应
%% method: 无
%% score: 4
%% degree: 850
%% duplicateId: 0
%% body:
在下列四个核反应方程中，x 表示质子的是\xzanswer{C}

\fourchoices[answer=C]
{$\ce{^{30}_{15}P -> ^{30}_{14}Si + x}$}
{$\ce{^{238}_{92}U -> ^{234}_{90}Th + x}$}
{$\ce{^{27}_{13}Al + ^{1}_{0}n -> ^{27}_{12}Mg + x}$}
{$\ce{^{27}_{13}Al + ^{4}_{2}He -> ^{30}_{15}P + x}$}

%% answer: C
%% memo:
\memoanswer{由核反应方程的质量数和电荷数守恒，可得各个选项中的 x 分别为正电子、$\alpha$ 粒子、质子、中子，故 C 正确。}

\item
%% number: 7
%% paperName: 2008年上海高考第 7 题 4 分
%% typeId: 1
%% type: 单选题
%% chapter: 力物体的平衡
%% point: 力矩平衡
%% method: 无
%% score: 4
%% degree: 650
%% duplicateId: 0
%% body:
如图所示，一根木棒 AB 在 O 点被悬挂起来，$AO=OC$，在 A、C 两点分别挂有二个和三个砝码，木棒处于平衡状态。如在木棒的 A、C 点各增加一个同样的砝码，则木棒\xzanswer{D}

\onepicture[w=7.0cm]{\includesvg{figs/07}}

\fourchoices[answer=D]
{绕 O 点顺时针方向转动}
{绕 O 点逆时针方向转动}
{平衡可能被破坏，转动方向不定}
{仍能保持平衡状态}

%% answer: D
%% memo:
\memoanswer{设木棒 AO 段重力为 $G_{1}$、重心离 O 点 $L_{1}$，BO 段重力为 $G_{2}$、重心离 O 点 $L_{2}$，$AO=OC=l$，由力矩平衡条件得 $G_{1}L_{1}+2Gl=G_{2}L_{2}+3Gl$；当两边各挂一个同样的砝码后，等式变为 $G_{1}L_{1}+3Gl=G_{2}L_{2}+4Gl$，依然成立，故木棒仍能保持平衡，D 正确。}

\item
%% number: 8
%% paperName: 2008年上海高考第 8 题 4 分
%% typeId: 1
%% type: 单选题
%% chapter: 功和能
%% point: 机械能守恒定律
%% method: 图象法
%% score: 4
%% degree: 450
%% duplicateId: 0
%% body:
物体做自由落体，$E_{k}$ 代表动能，$E_{p}$ 代表势能，$h$ 代表下落的距离，以水平地面为零势能面，下列所示图像中，能正确反映各物理量之间关系的是\xzanswer{B}

\onepicture[w=15.0cm]{\includesvg{figs/08}}

%% answer: B
%% memo:
\memoanswer{由机械能守恒定律得 $E_{p}=E-E_{k}$，故势能与动能的图像为倾斜的直线，C 错误；由动能定理得 $E_{k}=mgh=\frac{1}{2}mv^{2}=\frac{1}{2}mg^{2}t^{2}$，则 $E_{p}=E-mgh$，故势能与 $h$ 的图像也为倾斜的直线，D 错误；且 $E_{p}=E-\frac{1}{2}mv^{2}$，故势能与速度的图像为开口向下的抛物线，B 正确；同理 $E_{p}=E-\frac{1}{2}mg^{2}t^{2}$，势能与时间的图像也为开口向下的抛物线，A 错误。}

\item
%% number: 9
%% paperName: 2008年上海高考第 9 题 4 分
%% typeId: 1
%% type: 单选题
%% chapter: 气体性质
%% point: 气体图象
%% method: 无
%% score: 4
%% degree: 650
%% duplicateId: 0
%% body:
已知理想气体的内能与温度成正比。如图所示的实线为汽缸内一定质量的理想气体由状态 1 到状态 2 的变化曲线，则在整个过程中汽缸内气体的内能\xzanswer{B}

\onepicture[w=4.5cm]{\includesvg{figs/09}}

\fourchoices[answer=B]
{先增大后减小}
{先减小后增大}
{单调变化}
{保持不变}

%% answer: B
%% memo:
\memoanswer{由 $\frac{pV}{T}$ 为恒量，用图像与坐标轴围成的面积表示 $pV$ 乘积，从实线与虚线等温线比较可得出，该面积先减小后增大，说明温度 $T$ 先减小后增大，内能先减小后增大。}

\item
%% number: 10
%% paperName: 2008年上海高考第 10 题 4 分
%% typeId: 1
%% type: 单选题
%% chapter: 电磁感应
%% point: 切割磁感线电动势
%% method: 图象法
%% score: 4
%% degree: 450
%% duplicateId: 0
%% body:
如图所示，平行于 $y$ 轴的导体棒以速度 $v$ 向右做匀速运动，经过半径为 $R$、磁感应强度为 $B$ 的圆形匀强磁场区域，导体棒中的感应电动势 $E$ 与导体棒的位置 $x$ 关系的图像是\xzanswer{A}

\onepicture[w=4.5cm]{\includesvg{figs/10a}}

\onepicture[w=15.0cm]{\includesvg{figs/10b}}

%% answer: A
%% memo:
\memoanswer{在 $x=R$ 左侧，导体棒切割有效长度 $l=2\sqrt{R^{2}-(R-x)^{2}}$，电动势与有效长度成正比，故在 $x=R$ 左侧，电动势与 $x$ 的关系为椭圆图像关系；由对称性可知在 $x=R$ 右侧与左侧的图像对称，故 A 正确。}

\item
%% number: 11
%% paperName: 2008年上海高考第 11 题 5 分
%% typeId: 2
%% type: 多选题
%% chapter: 直线运动
%% point: 竖直上下抛运动
%% method: 无
%% score: 5
%% degree: 650
%% duplicateId: 0
%% body:
某物体以 $30\Ums$ 的初速度竖直上抛，不计空气阻力，$g$ 取 $10\Umsq$。5 s 内物体的\xzanswer{AB}

\fourchoices[answer=AB]
{路程为 $65\Um$}
{位移大小为 $25\Um$，方向向上}
{速度改变量的大小为 $10\Ums$}
{平均速度大小为 $13\Ums$，方向向上}

%% answer: AB
%% memo:
\memoanswer{初速度 $30\Ums$，只需 $3\Us$ 即可上升到最高点，上升高度为 $h_{1}=\frac{30^{2}}{20}\Um=45\Um$；再自由落体 $2\Us$，下降高度为 $h_{2}=\frac{1}{2}\times10\times2^{2}\Um=20\Um$。故总路程为 $65\Um$，A 正确；此时离抛出点高 $25\Um$，位移方向竖直向上，B 正确；此时速度为 $v=10\times2\Ums=20\Ums$，速度改变量为 $50\Ums$，C 错误；平均速度大小为 $\frac{25\Um}{5\Us}=5\Ums$，D 错误。}

\item
%% number: 12
%% paperName: 2008年上海高考第 12 题 5 分
%% typeId: 2
%% type: 多选题
%% chapter: 光学
%% point: 光的干涉
%% method: 无
%% score: 5
%% degree: 650
%% duplicateId: 0
%% body:
在杨氏双缝干涉实验中，如果\xzanswer{BD}

\fourchoices[answer=BD]
{用白光作为光源，屏上将呈现黑白相间的条纹}
{用红光作为光源，屏上将呈现红黑相间的条纹}
{用红光照射一条狭缝，用紫光照射另一条狭缝，屏上将呈现彩色条纹}
{用紫光作为光源，遮住其中一条狭缝，屏上将呈现间距不等的条纹}

%% answer: BD
%% memo:
\memoanswer{白光做杨氏双缝干涉实验，屏上将呈现彩色条纹，A 错误；用红光作光源，屏上将呈现红色亮条纹与暗条纹（即黑条纹）相间，B 正确；红光和紫光频率不同，不能产生干涉条纹，C 错误；紫光作光源，遮住一条狭缝，屏上出现单缝衍射条纹，即间距不等的条纹，D 正确。}

\item
%% number: 13
%% paperName: 2008年上海高考第 13 题 5 分
%% typeId: 2
%% type: 多选题
%% chapter: 气体性质
%% point: 液柱移动定性分析
%% method: 无
%% score: 5
%% degree: 650
%% duplicateId: 0
%% body:
如图所示，两端开口的弯管，左管插入水银槽中，右管有一段高为 $h$ 的水银柱，中间封有一段空气。则\xzanswer{ACD}

\onepicture[w=4.5cm]{\includesvg{figs/13}}

\fourchoices[answer=ACD]
{弯管左管内外水银面的高度差为 $h$}
{若把弯管向上移动少许，则管内气体体积增大}
{若把弯管向下移动少许，右管内的水银柱沿管壁上升}
{若环境温度升高，右管内的水银柱沿管壁上升}

%% answer: ACD
%% memo:
\memoanswer{封闭气体的压强等于大气压与水银柱产生压强之差，故左管内外水银面高度差也为 $h$，A 正确；弯管上下移动，封闭气体温度和压强不变，体积不变，B 错误，C 正确；环境温度升高，封闭气体体积增大，则右管内的水银柱沿管壁上升，D 正确。}

\item
%% number: 14
%% paperName: 2008年上海高考第 14 题 5 分
%% typeId: 2
%% type: 多选题
%% chapter: 电场
%% point: 带电粒子的往复运动
%% method: 无
%% score: 5
%% degree: 250
%% duplicateId: 0
%% body:
如图所示，在光滑绝缘水平面上，两个带等量正电的点电荷 M、N 分别固定在 A、B 两点，O 为 AB 连线的中点，CD 为 AB 的垂直平分线。在 CD 之间的 F 点由静止释放一个带负电的小球 P（设不改变原来的电场分布），在以后的一段时间内，P 在 CD 连线上做往复运动，则\xzanswer{BCD}

\onepicture[w=7.0cm]{\includesvg{figs/14}}

\fourchoices[answer=BCD]
{小球 P 的带电量缓慢减小，则它往复运动过程中的振幅不断减小}
{小球 P 的带电量缓慢减小，则它往复运动过程中每次经过 O 点时的速率不断减小}
{点电荷 M、N 的带电量同时等量地缓慢增大，则小球 P 往复运动过程中周期不断减小}
{点电荷 M、N 的带电量同时等量地缓慢增大，则小球 P 往复运动过程中的振幅不断减小}

%% answer: BCD
%% memo:
\memoanswer{设 F 与 $F'$ 绕 O 点对称，在 F 与 $F'$ 之间，小球始终受到指向 O 点的回复力并做往复运动。若小球 P 带电量缓慢减小，则此后小球能运动到 $F'$ 点下方，即振幅会加大，A 错误；每次经过 O 点因电场力做功减少而速率不断减小，B 正确；若点电荷 M、N 电荷量缓慢增大，则中垂线 CD 上的场强相对增大，振幅减小，加速度相对原来每个位置增大，故一个周期的时间必定减小，C、D 正确。}

\item
%% number: 15
%% paperName: 2008年上海高考第 15 题 4 分
%% typeId: 4
%% type: 实验题
%% chapter: 光学
%% point: 光电效应
%% method: 无
%% score: 4
%% degree: 650
%% duplicateId: 0
%% body:
如图所示，用导线将验电器与洁净锌板连接，触摸锌板使验电器指示归零，用紫外线照射锌板，验电器指针发生明显偏转，接着用毛皮摩擦过的橡胶棒接触锌板，发现验电器指针张角减小，此现象说明锌板带 \tkanswer[1.4cm]{正} 电（选填“正”或“负”）；若改用红外线重复以上实验，结果发现验电器指针根本不会偏转，说明金属锌的极限频率 \tkanswer[1.6cm]{大于} 红外线的频率（选填“大于”或“小于”）。

\onepicture[w=5.5cm, align=right]{\includesvg{figs/15}}

%% answer:
\jdanswer{正，大于}

%% memo:
\memoanswer{毛皮摩擦过的橡胶棒带负电，因锌板被紫外线照射后发生光电效应缺少电子而带正电，故验电器指针的负电荷与锌板正电荷中和一部分后偏角变小；用红外线照射验电器指针偏角不变，说明锌板未发生光电效应，锌板的极限频率大于红外线的频率。}

\item
%% number: 16
%% paperName: 2008年上海高考第 16 题 4 分
%% typeId: 4
%% type: 实验题
%% chapter: 光学
%% point: 光的干涉
%% method: 无
%% score: 4
%% degree: 650
%% duplicateId: 0
%% body:
用如图所示的实验装置观察光的薄膜干涉现象，图（a）是点燃的酒精灯（在灯芯上洒些盐），图（b）是竖立的附着一层肥皂液薄膜的金属线圈，将金属线圈在其所在的平面内缓慢旋转，观察到的现象是（　）

\onepicture[w=8.0cm]{\includesvg{figs/16}}

\fourchoices[answer=D]
{当金属线圈旋转 $30^{\circ}$ 时，干涉条纹同方向旋转 $30^{\circ}$}
{当金属线圈旋转 $45^{\circ}$ 时，干涉条纹同方向旋转 $90^{\circ}$}
{当金属线圈旋转 $60^{\circ}$ 时，干涉条纹同方向旋转 $30^{\circ}$}
{干涉条纹保持不变}

%% answer:
\jdanswer{D}

%% memo:
\memoanswer{金属丝圈的转动，改变不了肥皂液膜的上薄下厚的形状，由干涉原理可知干涉条纹与金属丝圈在该竖直平面内的转动无关，仍然是水平的干涉条纹，故 D 正确。}

\item
%% number: 17
%% paperName: 2008年上海高考第 17 题 6 分
%% typeId: 4
%% type: 实验题
%% chapter: 机械振动
%% point: 单摆测重力加速度
%% method: 无
%% score: 6
%% degree: 450
%% duplicateId: 0
%% body:
在“用单摆测重力加速度”的实验中。

\begin{enumerate}
	\item
	某同学的操作步骤为：
	
	A．取一根细线，下端系住直径为 $d$ 的金属小球，上端固定在铁架台上；
	
	B．用米尺量得细线长度 $l$；
	
	C．在细线偏离竖直方向 $5^{\circ}$ 位置释放小球；
	
	D．用秒表记录小球完成 $n$ 次全振动所用的总时间 $t$，得到周期 $T=t/n$；
	
	E．用公式 $g=\frac{4\pi^{2}l}{T^{2}}$ 计算重力加速度。
	
	按上述方法得出的重力加速度与实际值相比 \tkanswer[1.4cm]{偏小}（选填“偏大”、“相同”或“偏小”）。
	
	\item
	已知单摆在任意摆角 $\theta$ 时的周期公式可近似为 $T'=T_{0}\left[1+a\sin^{2}\frac{\theta}{2}\right]$，式中 $T_{0}$ 为摆角 $\theta$ 趋近于 $0^{\circ}$ 时的周期，$a$ 为常数。为了用图像法验证该关系式，需要测量的物理量有 \tkanswer[2.6cm]{$T'$、$\theta$}；若某同学在实验中得到了如图所示的图线，则图像中的横轴表示 \tkanswer[1.6cm]{$T'$}。
\end{enumerate}

\onepicture[w=5.0cm, align=right]{\includesvg{figs/17}}

%% answer:
\jdanswer{
\begin{enumerate}
	\item
	偏小
	
	\item
	$T'$（或 $t$、$n$）、$\theta$；$T'$
\end{enumerate}
}

%% memo:
\memoanswer{
（1）单摆的摆长应为摆线长度与小球半径之和，该同学将偏小的摆长代入公式计算，所得重力加速度的测量值偏小。

（2）为验证该关系式，需要测量单摆在任意摆角 $\theta$ 时的周期 $T'$，还需要测量摆角 $\theta$；由公式与图像的函数关系式可推导得到摆角 $\theta$ 趋近于 $0^{\circ}$ 时横轴的截距为 $T_{0}$，故图像中的横轴表示 $T'$。
}

\item
%% number: 18
%% paperName: 2008年上海高考第 18 题 6 分
%% typeId: 4
%% type: 实验题
%% chapter: 电路
%% point: 电路实验
%% method: 无
%% score: 6
%% degree: 450
%% duplicateId: 0
%% body:
某同学利用图（a）所示的电路研究灯泡 $L_{1}$（$6\UV$，$1.5\UW$）、$L_{2}$（$6\UV$，$10\UW$）的发光情况（假设灯泡电阻恒定），图（b）为实物图。

\begin{enumerate}
	\item
	他分别将 $L_{1}$、$L_{2}$ 接入图（a）中的虚线框位置，移动滑动变阻器的滑片，当电压表示数为 $6\UV$ 时，发现灯泡均能正常发光。在图（b）中用笔线代替导线将电路连线补充完整。
	
	\item
	接着他将 $L_{1}$ 和 $L_{2}$ 串联后接入图（a）中的虚线框位置，移动滑动变阻器的滑片，当电压表示数为 $6\UV$ 时，发现其中一个灯泡亮而另一个灯泡不亮，出现这种现象的原因是 \tkanswer[5.0cm]{}。
	
	\item
	现有如下器材：电源 $E$（$6\UV$，内阻不计），灯泡 $L_{1}$（$6\UV$，$1.5\UW$）、$L_{2}$（$6\UV$，$10\UW$）、$L_{3}$（$6\UV$，$10\UW$），单刀双掷开关 S，在图（c）中设计一个机动车转向灯的控制电路：当单刀双掷开关 S 与 1 相接时，信号灯 $L_{1}$ 亮，右转向灯 $L_{2}$ 亮而左转向灯 $L_{3}$ 不亮；当单刀双掷开关 S 与 2 相接时，信号灯 $L_{1}$ 亮，左转向灯 $L_{3}$ 亮而右转向灯 $L_{2}$ 不亮。
\end{enumerate}

\threepicture[labelA=2008上海18a, labelB=2008上海18b, labelC=2008上海18c, align=right]
{\includesvg{figs/18a}}
{\includesvg{figs/18b}}
{\includesvg{figs/18c}}

%% answer:
\jdanswer{
\begin{enumerate}
	\item
	如图 b。
	\item
	由于 $L_{1}$ 的电阻比 $L_{2}$ 大得多，所以 $L_{2}$ 两端分到的电压很小，虽然有电流流过，但功率很小，不能发光。
	\item
	如图 c。
\end{enumerate}
}

%% memo:
\memoanswer{
（1）将 $L_{1}$（或 $L_{2}$）与滑动变阻器、电流表、电源、开关串联，电压表并联在灯泡两端，电路连线补充完整如图 b 所示。

\onepicture[w=4.5cm]{\includesvg{figs/18_an1}}

（2）由于灯泡 $L_{2}$ 和 $L_{1}$ 额定电压相同，灯泡 $L_{2}$ 功率大得多，故 $R_{L_{2}}$ 比 $R_{L_{1}}$ 小得多，灯泡 $L_{2}$ 分得的电压很小，虽然有电流流过，但功率很小，不能发光。

（3）当单刀双掷开关 S 与 1 相接时，$L_{1}$ 与 $L_{2}$ 并联接在电源两端，$L_{3}$ 所在支路断开，故 $L_{1}$、$L_{2}$ 亮而 $L_{3}$ 不亮；S 与 2 相接时，$L_{1}$ 与 $L_{3}$ 并联接在电源两端，$L_{2}$ 所在支路断开，故 $L_{1}$、$L_{3}$ 亮而 $L_{2}$ 不亮，电路设计如图 c 所示。

\onepicture[w=4.5cm]{\includesvg{figs/18_an2}}
}

\item
%% number: 19
%% paperName: 2008年上海高考第 19 题 10 分
%% typeId: 4
%% type: 实验题
%% chapter: 电磁感应
%% point: 电磁感应的电量
%% method: 无
%% score: 10
%% degree: 450
%% duplicateId: 0
%% body:
如图所示是测量通电螺线管 A 内部磁感应强度 $B$ 及其与电流 $I$ 关系的实验装置。将截面积为 $S$、匝数为 $N$ 的小试测线圈 P 置于通电螺线管 A 中间，试测线圈平面与螺线管的轴线垂直，可认为穿过该试测线圈的磁场均匀，将试测线圈引线的两端与冲击电流计 D 相连。拨动双刀双掷换向开关 K，改变通入螺线管的电流方向，而不改变电流的大小，在 P 中产生的感应电流引起 D 的指针偏转。

\begin{enumerate}
	\item
	将开关合到位置 1，待螺线管中的电流稳定后，再将 K 从位置 1 拨到位置 2，测得 D 的最大偏转距离为 $d_{m}$，已知冲击电流计的磁通灵敏度为 $D_{\varphi}$，$D_{\varphi}=\frac{d_{m}}{N\Delta\varphi}$，式中 $\Delta\varphi$ 为单匝试测线圈磁通量的变化量，则试测线圈所在处的磁感应强度的大小为 $B=$ \tkanswer[3.0cm]{}；若将 K 从位置 1 拨到位置 2 所用的时间为 $\Delta t$，则试测线圈 P 中产生的平均感应电动势 $\varepsilon=$ \tkanswer[3.0cm]{}。
	
	\item
	调节可变电阻 $R$，多次改变电流并拨动 K，得到 A 中电流 $I$ 和磁感应强度 $B$ 的数据，见下表。由此可得，螺线管 A 内磁感应强度 $B$ 与电流 $I$ 的关系式为 $B=$ \tkanswer[2.6cm]{}。
	
	\begin{table}[htbp]
	\centering
	\begin{tblr}{|c|c|c|}
	\hline
	实验次数 & 电流 $I$（$\UA$） & 磁感应强度 $B$（$\times10^{-3}\UT$） \\
	\hline
	1 & 0.5 & 0.62 \\
	\hline
	2 & 1.0 & 1.25 \\
	\hline
	3 & 1.5 & 1.88 \\
	\hline
	4 & 2.0 & 2.51 \\
	\hline
	5 & 2.5 & 3.12 \\
	\hline
	\end{tblr}
	\end{table}
	
	\item
	（多选题）为了减少实验误差，提高测量的准确性，可采取的措施有（　）
	
	（A）适当增加试测线圈的匝数 $N$
	
	（B）适当增大试测线圈的横截面积 $S$
	
	（C）适当增大可变电阻 $R$ 的阻值
	
	（D）适当延长拨动开关的时间
\end{enumerate}

\onepicture[w=8.0cm, align=right]{\includesvg{figs/19}}

%% answer:
\jdanswer{
\begin{enumerate}
	\item
	$\frac{d_{m}}{2NSD_{\varphi}}$，$\frac{d_{m}}{D_{\varphi}\Delta t}$
	\item
	$B=1.25\times10^{-3}I$（或 $B=kI$）
	\item
	A、B
\end{enumerate}
}

%% memo:
\memoanswer{
（1）改变电流方向，磁通量变化量为原来磁通量的两倍，即 $2BS$，代入公式 $D_{\varphi}=\frac{d_{m}}{N\Delta\varphi}$ 计算得 $B=\frac{d_{m}}{2ND_{\varphi}S}$；由法拉第电磁感应定律可知平均感应电动势为 $\varepsilon=\frac{d_{m}}{D_{\varphi}\Delta t}$。

（2）由数据可得 $B$ 与 $I$ 成正比，比例常数约为 $1.25\times10^{-3}$，故 $B=1.25\times10^{-3}I$（或 $B=kI$）。

（3）为了得到平均电动势的准确值，时间要尽量小；由 $B$ 的计算值可看出与 $N$ 和 $S$ 相关联，故选择 A、B。
}

\item
%% number: 20
%% paperName: 2008年上海高考第 20 题 10 分
%% typeId: 6
%% type: 计算题
%% chapter: 交流电
%% point: 变压器
%% method: 无
%% score: 10
%% degree: 650
%% duplicateId: 0
%% body:
\textbf{A．}汽车行驶时轮胎的胎压太高容易造成爆胎事故，太低又会造成耗油量上升。已知某型号轮胎能在 $-40^{\circ}\mathrm{C}\sim90^{\circ}\mathrm{C}$ 正常工作，为使轮胎在此温度范围内工作时的最高胎压不超过 $3.5\Uatm$，最低胎压不低于 $1.6\Uatm$，那么在 $t=20^{\circ}\mathrm{C}$ 时给该轮胎充气，充气后的胎压在什么范围内比较合适？（设轮胎的体积不变）

\textbf{B．}某小型水电站输出功率为 $20\UkW$，输电线路总电阻是 $6\UO$。

\begin{enumerate}
	\item
	若采用 $380\UV$ 输电，求输电线路损耗的功率；
	\item
	若改用 $5000\UV$ 高压输电，用户端利用 $n_{1}:n_{2}=22:1$ 的变压器降压，求用户得到的电压。
\end{enumerate}

%% answer:
\jdanswer{
\textbf{A．}充气后的胎压在 $2.01\Uatm\sim2.83\Uatm$ 范围内比较合适。

\textbf{B．}
\begin{enumerate}
	\item
	$P_{\text{损}}=16.62\UkW$
	\item
	$U_{2}=226.18\UV$
\end{enumerate}
}

%% memo:
\memoanswer{
\textbf{A．}由于轮胎容积不变，轮胎内气体做等容变化。设在 $T_{0}=293\UK$ 充气后的最小胎压为 $p_{\text{min}}$、最大胎压为 $p_{\text{max}}$。依题意，当 $T_{1}=233\UK$ 时胎压为 $p_{1}=1.6\Uatm$，根据查理定律 $\frac{p_{1}}{T_{1}}=\frac{p_{\text{min}}}{T_{0}}$，即 $\frac{1.6}{233}=\frac{p_{\text{min}}}{293}$，解得 $p_{\text{min}}=2.01\Uatm$；当 $T_{2}=363\UK$ 时胎压为 $p_{2}=3.5\Uatm$，根据查理定律 $\frac{p_{2}}{T_{2}}=\frac{p_{\text{max}}}{T_{0}}$，即 $\frac{3.5}{363}=\frac{p_{\text{max}}}{293}$，解得 $p_{\text{max}}=2.83\Uatm$。故充气后的胎压在 $2.01\Uatm\sim2.83\Uatm$ 范围内比较合适。

\textbf{B．}（1）输电线上的电流强度为 $I=\frac{P}{U}=\frac{20\times10^{3}}{380}\UA=52.63\UA$，输电线损耗的功率为 $P_{\text{损}}=I^{2}R=52.63^{2}\times6\UW=16620\UW=16.62\UkW$。

（2）改用高压输电后，输电线上的电流强度变为 $I'=\frac{P}{U'}=\frac{20\times10^{3}}{5000}\UA=4\UA$；用户端在变压器降压前获得的电压 $U_{1}=U-I'R=(5000-4\times6)\UV=4976\UV$；根据 $\frac{U_{1}}{U_{2}}=\frac{n_{1}}{n_{2}}$，用户得到的电压为 $U_{2}=\frac{n_{2}}{n_{1}}U_{1}=\frac{1}{22}\times4976\UV=226.18\UV$。
}

\item
%% number: 21
%% paperName: 2008年上海高考第 21 题 12 分
%% typeId: 6
%% type: 计算题
%% chapter: 功和能
%% point: 动能定理
%% method: 无
%% score: 12
%% degree: 450
%% duplicateId: 0
%% body:
总质量为 $80\Ukg$ 的跳伞运动员从离地 $500\Um$ 的直升机上跳下，经过 $2\Us$ 拉开绳索开启降落伞，如图所示是跳伞过程中的 $v-t$ 图，试根据图像求：（$g$ 取 $10\Umsq$）

\begin{enumerate}
	\item
	$t=1\Us$ 时运动员的加速度和所受阻力的大小；
	\item
	估算 $14\Us$ 内运动员下落的高度及克服阻力做的功；
	\item
	估算运动员从飞机上跳下到着地的总时间。
\end{enumerate}

\onepicture[w=7.0cm, align=right]{\includesvg{figs/21}}

%% answer:
\jdanswer{
\begin{enumerate}
	\item
	$a=8\Umsq$，$f=160\UN$
	\item
	$h\approx158\Um$，$W_{f}\approx1.25\times10^{5}\UJ$
	\item
	$t_{\text{总}}=71\Us$
\end{enumerate}
}

%% memo:
\memoanswer{
（1）由图中可以看出，在 $t=2\Us$ 内运动员做匀加速运动，其加速度大小为 $a=\frac{v_{t}}{t}=\frac{16}{2}\Umsq=8\Umsq$。设此过程中运动员受到的阻力大小为 $f$，根据牛顿第二定律有 $mg-f=ma$，得 $f=mg-ma=80\times(10-8)\UN=160\UN$。

（2）从图中可以估算出运动员在 $14\Us$ 内下落了 $39.5\times2\times2\Um=158\Um$。根据动能定理有 $mgh-W_{f}=\frac{1}{2}mv^{2}$，所以 $W_{f}=mgh-\frac{1}{2}mv^{2}=(80\times10\times158-0.5\times80\times6^{2})\UJ\approx1.25\times10^{5}\UJ$。

（3）$14\Us$ 后运动员做匀速运动的时间为 $t'=\frac{H-h}{v_{t}}=\frac{500-158}{6}\Us=57\Us$，运动员从飞机上跳下到着地需要的总时间 $t_{\text{总}}=t+t'=(14+57)\Us=71\Us$。
}

\item
%% number: 22
%% paperName: 2008年上海高考第 22 题 12 分
%% typeId: 5
%% type: 辨析题
%% chapter: 机械波
%% point: 波的多解性
%% method: 无
%% score: 12
%% degree: 450
%% duplicateId: 0
%% body:
有两列简谐横波 a、b 在同一媒质中沿 $x$ 轴正方向传播，波速均为 $v=2.5\Ums$。在 $t=0$ 时两列波的波峰正好在 $x=2.5\Um$ 处重合，如图所示。

\begin{enumerate}
	\item
	求两列波的周期 $T_{a}$ 和 $T_{b}$；
	\item
	求 $t=0$ 时两列波的波峰重合处的所有位置；
	\item
	辨析题：分析和判断在 $t=0$ 时是否存在两列波的波谷重合处。某同学分析如下：既然两列波的波峰与波峰存在重合处，那么波谷与波谷重合处也一定存在。只要找到这两列波半波长的最小公倍数，……，即可得到波谷与波谷重合处的所有位置。你认为该同学的分析正确吗？若正确，求出这些位置；若不正确，指出错误处并通过计算说明理由。
\end{enumerate}

\onepicture[w=14.0cm, align=right]{\includesvg{figs/22}}

%% answer:
\jdanswer{
\begin{enumerate}
	\item
	$T_{a}=1\Us$，$T_{b}=1.6\Us$
	\item
	$x=2.5\pm20k\Um$，$k=0,1,2,3\cdots$
	\item
	该同学的分析不正确。要找两列波的波谷与波谷重合处，必须从波峰重合处出发，找到这两列波半波长的奇数倍恰好相等的位置。设距离 $x=2.5\Um$ 为 $L$ 处两列波的波谷与波谷相遇，并设 $L=(2m-1)\frac{\lambda_{a}}{2}$，$L=(2n-1)\frac{\lambda_{b}}{2}$（$m$、$n$ 均为正整数）。只要能找到对应的 $m$、$n$ 即可。将 $\lambda_{a}=2.5\Um$、$\lambda_{b}=4\Um$ 代入并整理，得 $\frac{2m-1}{2n-1}=\frac{\lambda_{b}}{\lambda_{a}}=\frac{4}{2.5}=\frac{8}{5}$，由于上式中的 $m$、$n$ 在正整数范围内无解，所以不存在波谷与波谷重合处。
\end{enumerate}
}

%% memo:
\memoanswer{
（1）从图中可以看出两列波的波长分别为 $\lambda_{a}=2.5\Um$、$\lambda_{b}=4\Um$，因此它们的周期分别为 $T_{a}=\frac{\lambda_{a}}{v}=\frac{2.5}{2.5}\Us=1\Us$，$T_{b}=\frac{\lambda_{b}}{v}=\frac{4.0}{2.5}\Us=1.6\Us$。

（2）两列波长的最小公倍数为 $s=20\Um$，在 $t=0$ 时，两列波的波峰重合处的所有位置为 $x=2.5\pm20k\Um$，$k=0,1,2,3\cdots$。

（3）该同学的分析不正确。要找两列波的波谷与波谷重合处，必须从波峰重合处出发，找到这两列波半波长的奇数倍恰好相等的位置。设距离 $x=2.5\Um$ 为 $L$ 处两列波的波谷与波谷相遇，并设 $L=(2m-1)\frac{\lambda_{a}}{2}$，$L=(2n-1)\frac{\lambda_{b}}{2}$（$m$、$n$ 均为正整数），将 $\lambda_{a}=2.5\Um$、$\lambda_{b}=4\Um$ 代入整理得 $\frac{2m-1}{2n-1}=\frac{\lambda_{b}}{\lambda_{a}}=\frac{4}{2.5}=\frac{8}{5}$，即 $5(2m-1)=8(2n-1)$，左边为奇数、右边为偶数，矛盾，故 $m$、$n$ 在正整数范围内无解，不存在波谷与波谷重合处。
}

\item
%% number: 23
%% paperName: 2008年上海高考第 23 题 12 分
%% typeId: 6
%% type: 计算题
%% chapter: 电场
%% point: 带电粒子的偏转
%% method: 无
%% score: 12
%% degree: 450
%% duplicateId: 0
%% body:
如图所示为研究电子枪中电子在电场中运动的简化模型示意图。在 $Oxy$ 平面的 ABCD 区域内，存在两个大小均为 $E$ 的匀强电场 Ⅰ 和 Ⅱ，两电场的边界均是边长为 $L$ 的正方形（不计粒子所受重力）。

\begin{enumerate}
	\item
	在该区域 AB 边的中点处由静止释放电子，求电子离开 ABCD 区域的位置；
	\item
	在电场 Ⅰ 区域内适当位置由静止释放电子，电子恰能从 ABCD 区域左下角 D 处离开，求所有释放点的位置；
	\item
	若将左侧电场 Ⅱ 整体水平向右移动 $\frac{L}{n}$（$n\geqslant1$），仍使电子从 ABCD 区域左下角 D 处离开（D 不随电场移动），求在电场 Ⅰ 区域内由静止释放电子的所有位置。
\end{enumerate}

\onepicture[w=9.0cm, align=right]{\includesvg{figs/23}}

%% answer:
\jdanswer{
\begin{enumerate}
	\item
	$(-2L,\frac{1}{4}L)$
	\item
	$xy=\frac{L^{2}}{4}$
	\item
	$xy=L^{2}\left(\frac{1}{2n}+\frac{1}{4}\right)$
\end{enumerate}
}

%% memo:
\memoanswer{
（1）设电子的质量为 $m$、电量为 $q$，电子在电场 Ⅰ 中做匀加速直线运动，出区域 Ⅰ 时的速度为 $v_{0}$，此后进入电场 Ⅱ 做类平抛运动。假设电子从 CD 边射出，出射点纵坐标为 $y$，有 $eEL=\frac{1}{2}mv_{0}^{2}$，$\frac{L}{2}-y=\frac{1}{2}at^{2}=\frac{eE}{2m}\times\frac{L^{2}}{v_{0}^{2}}$，解得 $y=\frac{1}{4}L$，所以原假设成立，即电子离开 ABCD 区域的位置坐标为 $(-2L,\frac{1}{4}L)$。

（2）设释放点在电场区域 Ⅰ 中，其坐标为 $(x,y)$，在电场 Ⅰ 中电子被加速到 $v_{1}$，然后进入电场 Ⅱ 做类平抛运动，并从 D 点离开，有 $eEx=\frac{1}{2}mv_{1}^{2}$，$y=\frac{1}{2}at^{2}=\frac{eE}{2m}\times\frac{L^{2}}{v_{1}^{2}}$，解得 $xy=\frac{L^{2}}{4}$，即在电场 Ⅰ 区域内满足该方程的点即为所求位置。

（3）设电子从 $(x,y)$ 点释放，在电场 Ⅰ 中加速到 $v_{2}$，进入电场 Ⅱ 做类平抛运动，在高度为 $y'$ 处离开电场 Ⅱ 时的情景与（2）中类似，然后电子做匀速直线运动，经过 D 点，则有 $eEx=\frac{1}{2}mv_{2}^{2}$，$y-y'=\frac{1}{2}at^{2}=\frac{eE}{2m}\times\frac{L^{2}}{v_{2}^{2}}$，$v_{y}=at=\frac{eEL}{mv_{2}}$，$y'=v_{y}\cdot\frac{L}{nv_{2}}$，解得 $xy=L^{2}\left(\frac{1}{2n}+\frac{1}{4}\right)$，即在电场 Ⅰ 区域内满足该方程的点即为所求位置。
}

\item
%% number: 24
%% paperName: 2008年上海高考第 24 题 14 分
%% typeId: 6
%% type: 计算题
%% chapter: 电磁感应
%% point: 电磁感应中的匀变速
%% method: 无
%% score: 14
%% degree: 450
%% duplicateId: 0
%% body:
如图所示，竖直平面内有一半径为 $r$、电阻为 $R_{1}$、粗细均匀的光滑半圆形金属环，在 M、N 处与距离为 $2r$、电阻不计的平行光滑金属导轨 ME、NF 相接，EF 之间接有电阻 $R_{2}$，已知 $R_{1}=12R$，$R_{2}=4R$。在 MN 上方及 CD 下方有水平方向的匀强磁场 Ⅰ 和 Ⅱ，磁感应强度大小均为 $B$。现有质量为 $m$、电阻不计的导体棒 ab，从半圆环的最高点 A 处由静止下落，在下落过程中导体棒始终保持水平，与半圆形金属环及轨道接触良好，设平行导轨足够长。已知导体棒下落 $\frac{r}{2}$ 时的速度大小为 $v_{1}$，下落到 MN 处时的速度大小为 $v_{2}$。

\begin{enumerate}
	\item
	求导体棒 ab 从 A 处下落 $\frac{r}{2}$ 时的加速度大小；
	\item
	若导体棒 ab 进入磁场 Ⅱ 后棒中电流大小始终不变，求磁场 Ⅰ 和 Ⅱ 之间的距离 $h$ 和 $R_{2}$ 上的电功率 $P_{2}$；
	\item
	若将磁场 Ⅱ 的 CD 边界略微下移，导体棒 ab 进入磁场 Ⅱ 时的速度大小为 $v_{3}$，要使其在外力 $F$ 作用下做匀加速直线运动，加速度大小为 $a$，求所加外力 $F$ 随时间变化的关系式。
\end{enumerate}

\onepicture[w=4.5cm, align=right]{\includesvg{figs/24}}

%% answer:
\jdanswer{
\begin{enumerate}
	\item
	$a=g-\frac{3B^{2}r^{2}v_{1}}{4Rm}$
	\item
	$h=\frac{9R^{2}gm^{2}}{32B^{4}r^{4}}-\frac{v_{2}^{2}}{2g}$，$P_{2}=\frac{9Rg^{2}m^{2}}{16B^{2}r^{2}}$
	\item
	$F=ma-mg+\frac{4B^{2}r^{2}}{3R}(v_{3}+at)$
\end{enumerate}
}

%% memo:
\memoanswer{
（1）以导体棒为研究对象，棒在磁场 Ⅰ 中切割磁感线，棒中产生感应电动势。导体棒 ab 从 A 下落 $\frac{r}{2}$ 时，导体棒在重力与安培力作用下做加速运动，由牛顿第二定律得 $mg-BIl=ma$，式中 $l=\sqrt{3}r$；又 $I=\frac{Blv_{1}}{R_{\text{总}}}$，式中 $R_{\text{总}}=\frac{64R^{2}}{8R+8R}=4R$。由以上各式得 $a=g-\frac{3B^{2}r^{2}v_{1}}{4Rm}$。

（2）当导体棒 ab 通过磁场 Ⅱ 时，若安培力恰好等于重力，棒中电流大小始终不变，即 $mg=BI\cdot2r=B\cdot\frac{B\cdot2r\cdot v_{t}}{R_{\text{并}}}\cdot2r=\frac{4B^{2}r^{2}v_{t}}{R_{\text{并}}}$，式中 $R_{\text{并}}=\frac{48R^{2}}{12R+4R}=3R$，解得 $v_{t}=\frac{3Rmg}{4B^{2}r^{2}}$。导体棒从 MN 到 CD 做加速度为 $g$ 的匀加速直线运动，有 $v_{t}^{2}-v_{2}^{2}=2gh$，得 $h=\frac{9R^{2}gm^{2}}{32B^{4}r^{4}}-\frac{v_{2}^{2}}{2g}$。此时导体棒重力的功率为 $P_{G}=mgv_{t}=\frac{3m^{2}g^{2}R}{4B^{2}r^{2}}$，根据能量守恒定律，此时导体棒重力的功率完全转化为电路中的电功率，即 $P_{\text{电}}=P_{1}+P_{2}=P_{G}=\frac{3m^{2}g^{2}R}{4B^{2}r^{2}}$，所以 $P_{2}=\frac{3}{4}P_{G}=\frac{9Rg^{2}m^{2}}{16B^{2}r^{2}}$。

（3）设导体棒 ab 进入磁场 Ⅱ 后经过时间 $t$ 的速度大小为 $v_{t}'$，此时的安培力大小为 $F'=\frac{4B^{2}r^{2}v_{t}'}{3R}$。由于导体棒 ab 做匀加速直线运动，有 $v_{t}'=v_{3}+at$；根据牛顿第二定律有 $F+mg-F'=ma$，由以上各式解得 $F=\frac{4B^{2}r^{2}}{3R}(v_{3}+at)+ma-mg$。
}

\end{enumerate}

\end{document}
